\documentclass{beamer}
\usepackage{ctex}
\usepackage[T1]{fontenc}

% other packages
\usepackage{latexsym,xcolor,multicol,booktabs,calligra}
\usepackage{pstricks,listings,stackengine,cite}
\usepackage{anyfontsize}
\usepackage{hyperref}
\usepackage{graphicx}
\usepackage{amsmath,amsthm}

\setbeamertemplate{bibliography item}[text]
% \bibliographystyle{plain}
% 如果参考文献太多的话，可以像下面这样调整字体：
% \tiny\bibliographystyle{plain}

% \usefonttheme[onlymath]{serif}

% \newtheorem{theorem}{定理}[subsection]
% \newtheorem{lemma}{引理}[subsection]
% \newtheorem*{definition}{定义}
% \newtheorem{property}{性质}[subsection]
% \newtheorem{infer}{推论}[subsection]

\usepackage{color}

\newcommand{\bsym}[1]{\boldsymbol{#1}}
\newcommand{\mypar}[1]{\left( #1 \right)}
\newcommand{\gram}[1]{\bsym{G}\mypar{#1}}
\newcommand{\abs}[1]{\left|#1 \right|}
\newcommand{\Abs}[1]{\left\Vert#1\right\Vert}
\newcommand{\setof}[1]{\left\{#1 \right\}}
\newcommand{\indot}[2]{\left\langle #1,#2 \right\rangle}
\newcommand{\mat}[1]{\left[ #1 \right]}
\newcommand{\myvec}[1]{\left[ #1 \right]^\top}
\newcommand{\sqmat}[3]{\begin{#1}
{#2}_{11} & {#2}_{12} & \cdots & {#2}_{1{#3}} \\
{#2}_{21} & {#2}_{22} & \cdots & {#2}_{2{#3}} \\
\vdots & \vdots &   \ddots     & \vdots \\
{#2}_{{#3}1} & {#2}_{{#3}2} & \cdots & {#2}_{{#3}{#3}}
\end{#1}}
\newcommand{\normmat}[4]{\begin{#1}
{#2}_{11} & {#2}_{12} & \cdots & {#2}_{1{#4}} \\
{#2}_{21} & {#2}_{22} & \cdots & {#2}_{2{#4}} \\
\vdots & \vdots &   \ddots     & \vdots \\
{#2}_{{#3}1} & {#2}_{{#3}2} & \cdots & {#2}_{{#3}{#4}}
\end{#1}}
\newcommand{\nullsp}[1]{\bsym{\mathrm{N}}\mypar{#1}}
\newcommand{\colsp}[1]{\bsym{\mathrm{C}}\mypar{#1}}
\newcommand{\func}[2]{\mathrm{#1}\mypar{#2}}
\newcommand{\entry}[3]{\func{entry}{#1,#2,#3}}
\newcommand{\row}[2]{\func{row}{#1,#2}}
\newcommand{\col}[2]{\func{col}{#1,#2}}
\newcommand{\trace}[1]{\func{Tr}{#1}}
\newcommand{\diag}[1]{\func{diag}{#1}}
\newcommand{\rank}[1]{\func{rank}{#1}}
\newcommand{\adj}[1]{\func{adj}{#1}}
\newcommand{\myspan}[1]{\func{span}{#1}}
\newcommand{\enums}[2]{{#1}_1,{#1}_2,\dots,{#1}_{#2}}
\newcommand{\inv}[1]{{#1}^{-1}}
\newcommand{\pinv}[1]{\inv{\mypar{#1}}}
\newcommand{\ortcom}[1]{{#1}^{\bot}}

\renewcommand{\det}[1]{\func{det}{#1}}

\newcommand{\dif}{\mathrm{d}}
\newcommand{\fracdif}[1]{\frac{\dif}{\dif #1}}

\newcommand{\todo}[1]{{ \textcolor{red}{ TODO: #1}}}

\newcommand{\mata}{\bsym{A}}
\newcommand{\matb}{\bsym{B}}
\newcommand{\matc}{\bsym{C}}
\newcommand{\matd}{\bsym{D}}
\newcommand{\mate}{\bsym{E}}
\newcommand{\matf}{\bsym{F}}
\newcommand{\matfstar}{\bsym{F^*}}
\newcommand{\matg}{\bsym{G}}
\newcommand{\matH}{\bsym{H}}
\newcommand{\mati}{\bsym{I}}
\newcommand{\matj}{\bsym{J}}
\newcommand{\matk}{\bsym{K}}
\newcommand{\matl}{\bsym{L}}
\newcommand{\matlam}{\bsym{\Lambda}}
\newcommand{\matm}{\bsym{M}}
\newcommand{\matn}{\bsym{N}}
\newcommand{\mato}{\bsym{O}}
\newcommand{\matp}{\bsym{P}}
\newcommand{\matq}{\bsym{Q}}
\newcommand{\matr}{\bsym{R}}
\newcommand{\mats}{\bsym{S}}
\newcommand{\matsig}{\bsym{\Sigma}}
\newcommand{\matt}{\bsym{T}}
\newcommand{\matu}{\bsym{U}}
\newcommand{\matv}{\bsym{V}}
\newcommand{\matw}{\bsym{W}}
\newcommand{\matx}{\bsym{X}}
\newcommand{\maty}{\bsym{Y}}
\newcommand{\matz}{\bsym{Z}}

\newcommand{\field}{\bsym{\mathrm{F}}}
\newcommand{\rea}{\bsym{\mathbb{R}}}

\newcommand{\veca}{\bsym{a}}
\newcommand{\vecal}{\bsym{\alpha}}
\newcommand{\vecb}{\bsym{b}}
\newcommand{\vecbeta}{\bsym{\beta}}
\newcommand{\vecc}{\bsym{c}}
\newcommand{\vecd}{\bsym{d}}
\newcommand{\vecdelta}{\bsym{\delta}}
\newcommand{\vece}{\bsym{e}}
\newcommand{\veceps}{\bsym{\epsilon}}
\newcommand{\vecveps}{\bsym{\varepsilon}}
\newcommand{\vecf}{\bsym{f}}
\newcommand{\vecg}{\bsym{g}}
\newcommand{\vecgamma}{\bsym{\gamma}}
\newcommand{\vech}{\bsym{h}}
\newcommand{\veceta}{\bsym{\eta}}
\newcommand{\veci}{\bsym{i}}
\newcommand{\vecj}{\bsym{j}}
\newcommand{\veck}{\bsym{k}}
\newcommand{\vecl}{\bsym{l}}
\newcommand{\vecm}{\bsym{m}}
\newcommand{\vecn}{\bsym{n}}
\newcommand{\veco}{\bsym{o}}
\newcommand{\vecone}{\bsym{1}}
\newcommand{\vecp}{\bsym{p}}
\newcommand{\vecq}{\bsym{q}}
\newcommand{\vecr}{\bsym{r}}
\newcommand{\vecs}{\bsym{s}}
\newcommand{\vect}{\bsym{t}}
\newcommand{\vecu}{\bsym{u}}
\newcommand{\vecv}{\bsym{v}}
\newcommand{\vecw}{\bsym{w}}
\newcommand{\vecx}{\bsym{x}}
\newcommand{\vecxi}{\bsym{\xi}}
\newcommand{\vecy}{\bsym{y}}
\newcommand{\vecz}{\bsym{z}}
\newcommand{\veczero}{\bsym{0}}


\author{夏海淞}
\title{线性代数习题课}
\subtitle{第四章}
% \institute{School of Computer Science, Fudan University}
\date{\today}
\usepackage{College}

% defs
\def\cmd#1{\texttt{\color{red}\footnotesize $\backslash$#1}}
\def\env#1{\texttt{\color{blue}\footnotesize #1}}
\definecolor{deepblue}{rgb}{0,0,0.5}
\definecolor{deepred}{rgb}{0.6,0,0}
\definecolor{deepgreen}{rgb}{0,0.5,0}
\definecolor{halfgray}{gray}{0.55}

\lstset{
    basicstyle=\ttfamily\small,
    keywordstyle=\bfseries\color{deepblue},
    emphstyle=\ttfamily\color{deepred},    % Custom highlighting style
    stringstyle=\color{deepgreen},
    numbers=left,
    numberstyle=\small\color{halfgray},
    rulesepcolor=\color{red!20!green!20!blue!20},
    frame=shadowbox,
}

\begin{document}
\setlength{\parskip}{0.45\baselineskip}

\begin{frame}
    \titlepage
    \begin{figure}[htpb]
        \begin{center}
            \vspace*{-0.5cm}
            \includegraphics[width=0.1\linewidth]{../pic/FDU.jpeg}
        \end{center}
    \end{figure}
\end{frame}

\subsection*{4.1}
\begin{frame}
    \frametitle{题面}
    在次数不大于\(2\)的多项式线性空间\(P_2[x]\)中，试证：\(f_1=1\)，\(f_2=x-1\)，\(f_3=\mypar{x-2}\mypar{x-1}\)线性无关。
\end{frame}

\begin{frame}
    \frametitle{解答}
    由例4.2.2可知，\(\vecveps_1=1\)，\(\vecveps_2=x\)，\(\vecveps_3=x^2\)是\(P_2[x]\)的一组基。又因为
    \begin{align*}
        \begin{bmatrix}
            \vecf_1 & \vecf_2 & \vecf_3
        \end{bmatrix}
         & =
        \begin{bmatrix}
            \vecveps_1 & \vecveps_2 & \vecveps_3
        \end{bmatrix}
        \begin{bmatrix}
            1 & -1 & 2  \\
            0 & 1  & -3 \\
            0 & 0  & 1
        \end{bmatrix} \\
         & =
        \begin{bmatrix}
            \vecveps_1 & \vecveps_2 & \vecveps_3
        \end{bmatrix}\mata
    \end{align*}
    其中\(\rank{\mata}=3\)。由例4.2.4可知\(\vecf_1,\vecf_2,\vecf_3\)线性无关。
\end{frame}

\subsection*{4.2}
\begin{frame}
    \frametitle{题面}
    证明：以下三个多项式为\(P_2[x]\)的一组基：
    \begin{equation*}
        f_1=1,f_2=x-1,f_3=\mypar{x-1}^2
    \end{equation*}
    再求\(g(x)=5x^2+x+3\)在此基下的坐标。
\end{frame}

\begin{frame}
    \frametitle{解答}
    和习题4.1类似，令\(\vecveps_1=1\)，\(\vecveps_2=x\)，\(\vecveps_3=x^2\)表示\(P_2[x]\)的一组基。则有
    \begin{align*}
        \begin{bmatrix}
            \vecf_1 & \vecf_2 & \vecf_3
        \end{bmatrix}
         & =
        \begin{bmatrix}
            \vecveps_1 & \vecveps_2 & \vecveps_3
        \end{bmatrix}
        \begin{bmatrix}
            1 & -1 & 1  \\
            0 & 1  & -2 \\
            0 & 0  & 1
        \end{bmatrix} \\
         & =
        \begin{bmatrix}
            \vecveps_1 & \vecveps_2 & \vecveps_3
        \end{bmatrix}\mata
    \end{align*}
    其中\(\rank{\mata}=3\)，因此\(\mata\)可逆，\(\vecf_1,\vecf_2,\vecf_3\)线性无关。

    \pause

    对\(\forall\vecf\in P_2[x]\)，设其坐标为\(\vecx\)。则有
    \begin{equation*}
        \vecf=
        \begin{bmatrix}
            \vecveps_1 & \vecveps_2 & \vecveps_3
        \end{bmatrix}\vecx=
        \begin{bmatrix}
            \vecf_1 & \vecf_2 & \vecf_3
        \end{bmatrix}\inv{\mata}\vecx
    \end{equation*}
    因此\(P_2[x]\)中的任意向量可由\(\vecf_1,\vecf_2,\vecf_3\)线性表出。
    将\(g\mypar{x}\)在基\(\vecveps_1,\vecveps_2,\vecveps_3\)下的坐标代入即可得到\(g\mypar{x}\)在基\(\vecf_1,\vecf_2,\vecf_3\)下的坐标：
    \begin{equation*}
        \inv{\mata}\begin{bmatrix}3&1&5\end{bmatrix}^\top=\begin{bmatrix}9&11&5\end{bmatrix}^\top
    \end{equation*}

    综上，\(\vecf_1,\vecf_2,\vecf_3\)为\(P_2[x]\)的一组基，\(g\mypar{x}\)在此基下的坐标为\(\myvec{9,11,5}\)。
\end{frame}

\subsection*{4.3}
\begin{frame}
    \frametitle{题面}
    在次数不大于\(3\)的多项式空间\(P_3[x]\)中，
    \begin{enumerate}
        \item 求由基\(1,x,x^2,x^3\)到基\(1,1+x,\mypar{1+x}^2,\mypar{1+x}^3\)的过渡矩阵；
        \item 求\(f(x)=a_0+a_1x+a_2x^2+a_3x^3\)在基\(1,1+x,\mypar{1+x}^2,\mypar{1+x}^3\)下的坐标。
    \end{enumerate}
\end{frame}

\begin{frame}[allowframebreaks]
    \frametitle{解答}
    \begin{enumerate}
        \item {
              容易发现
              \begin{equation*}
                  \begin{bmatrix}
                      1 & 1+x & \mypar{1+x}^2 & \mypar{1+x}^3
                  \end{bmatrix}=
                  \begin{bmatrix}
                      1 \\ x \\ x^2 \\ x^3
                  \end{bmatrix}^\top
                  \begin{bmatrix}
                      1 & 1 & 1 & 1 \\
                      0 & 1 & 2 & 3 \\
                      0 & 0 & 1 & 3 \\
                      0 & 0 & 0 & 1
                  \end{bmatrix}
              \end{equation*}
              因此过渡矩阵为\(\matm=
              \begin{bmatrix}
                  1 & 1 & 1 & 1 \\
                  0 & 1 & 2 & 3 \\
                  0 & 0 & 1 & 3 \\
                  0 & 0 & 0 & 1
              \end{bmatrix}\)；
              }
        \item {
              令\(\veca=\myvec{a_0,a_1,a_2,a_3}\)，则有
              \begin{align*}
                  f\mypar{x}
                   & =
                  \begin{bmatrix}
                      1 & x & x^2 & x^3
                  \end{bmatrix}\veca \\
                   & =
                  \begin{bmatrix}
                      1 & 1+x & \mypar{1+x}^2 & \mypar{1+x}^3
                  \end{bmatrix}\inv{\matm}\veca
              \end{align*}
              因此\(f\mypar{x}\)在基\(1,1+x,\mypar{1+x}^2,\mypar{1+x}^3\)下的坐标为
              \begin{equation*}
                  \inv{\matm}\veca=
                  \begin{bmatrix}
                      a_0-a_1+a_2-a_3 \\
                      a_1-2a_2+3a_3   \\
                      a_2-3a_3        \\
                      a_3
                  \end{bmatrix}
              \end{equation*}
              }
    \end{enumerate}
\end{frame}

\subsection*{4.4}
\begin{frame}
    \frametitle{题面}
    在\(P_3[x]\)的多项式空间中，旧基为\(1,x,x^2,x^3\)；新基为\(1\)，\(1+x\)，\(1+x+x^2\)，\(1+x+x^2+x^3\)。
    \begin{enumerate}
        \item 求旧基到新基的过渡矩阵；
        \item 求多项式\(1+2x+3x^2+4x^3\)在新基下的坐标；
        \item 若多项式\(f(x)\)在新基下的坐标为\(\myvec{1,2,3,4}\)，求它在旧基下的坐标。
    \end{enumerate}
\end{frame}

\begin{frame}
    \frametitle{解答}
    \begin{enumerate}
        \item 过渡矩阵为\(\matm=\begin{bmatrix}1&1&1&1\\0&1&1&1\\0&0&1&1\\0&0&0&1\end{bmatrix}\)；\pause
        \item 因为该多项式在旧基下的坐标为\(\vecx=\myvec{1,2,3,4}\)，因此该多项式在新基下的坐标为\(\inv{\matm}\vecx=\myvec{-1,-1,-1,4}\)；\pause
        \item 同理可得该多项式在旧基下的坐标为\(\matm\vecx=\myvec{10,9,7,4}\)。
    \end{enumerate}
\end{frame}

\subsection*{4.5}
\begin{frame}
    \frametitle{题面}
    已知\(\vecxi\)在基\(\matb_1=\setof{\vecal_1,\vecal_2,\vecal_3}\)下的坐标为\(\vecxi_{\matb_1}=\myvec{1,-2,2}\)，求\(\vecxi\)在基\(\matb_2=\setof{\vecbeta_1,\vecbeta_2,\vecbeta_3}\)下的坐标\(\vecxi_{\matb_2}\)，其中\(\vecbeta_1=\vecal_1+\vecal_2\)，\(\vecbeta_2=\vecal_2+\vecal_3\)，\(\vecbeta_3=\vecal_3+\vecal_1\)。
\end{frame}

\begin{frame}
    \frametitle{解答}
    由题设可知
    \begin{equation*}
        \begin{bmatrix}\vecbeta_1&\vecbeta_2&\vecbeta_3\end{bmatrix}=
        \begin{bmatrix}\vecal_1&\vecal_2&\vecal_3\end{bmatrix}
        \begin{bmatrix}
            1 & 0 & 1 \\
            1 & 1 & 0 \\
            0 & 1 & 1
        \end{bmatrix}
    \end{equation*}
    因此从基\(\matb_1\)到基\(\matb_2\)的过渡矩阵为\(\matm=\begin{bmatrix}1&0&1\\1&1&0\\0&1&1\end{bmatrix}\)。
    因此有
    \begin{equation*}
        \vecxi_{\matb_2}=\inv{\matm}\vecxi_{\matb_1}=\begin{bmatrix}-1.5&-0.5&2.5\end{bmatrix}^\top
    \end{equation*}
\end{frame}

\subsection*{4.6}
\begin{frame}
    \frametitle{题面}
    设\(\vecveps_1,\vecveps_2,\vecveps_3\)是线性空间\(V\)的一组基，且
    \begin{equation*}
        \begin{cases}
            \vecxi_1=\vecveps_1+\vecveps_3 \\
            \vecxi_2=\vecveps_2            \\
            \vecxi_3=\vecveps_1+2\vecveps_2+2\vecveps_3
        \end{cases}
    \end{equation*}
    \begin{equation*}
        \begin{cases}
            \veceta_1=\vecveps_1            \\
            \veceta_2=\vecveps_1+\vecveps_2 \\
            \veceta_3=\vecveps_1+\vecveps_2+\vecveps_3
        \end{cases}
    \end{equation*}
    \begin{enumerate}
        \item 试证\(\vecxi_1,\vecxi_2,\vecxi_3\)及\(\veceta_1,\veceta_2,\veceta_3\)都是\(V\)的一组基；
        \item 求由基\(\vecxi_1,\vecxi_2,\vecxi_3\)到基\(\veceta_1,\veceta_2,\veceta_3\)的过渡矩阵。
    \end{enumerate}
\end{frame}

\begin{frame}[allowframebreaks]
    \frametitle{解答}
    \begin{enumerate}
        \item {
              容易发现
              \begin{align*}
                  \begin{bmatrix}\vecxi_1&\vecxi_2&\vecxi_3\end{bmatrix}
                   & =
                  \begin{bmatrix}\vecveps_1&\vecveps_2&\vecveps_3\end{bmatrix}
                  \begin{bmatrix}
                      1 & 0 & 1 \\
                      0 & 1 & 2 \\
                      1 & 0 & 2
                  \end{bmatrix}                                                        \\
                   & =\begin{bmatrix}\vecveps_1&\vecveps_2&\vecveps_3\end{bmatrix}\mata
              \end{align*}
              \begin{align*}
                  \begin{bmatrix}\veceta_1&\veceta_2&\veceta_3\end{bmatrix}
                   & =
                  \begin{bmatrix}\vecveps_1&\vecveps_2&\vecveps_3\end{bmatrix}
                  \begin{bmatrix}
                      1 & 1 & 1 \\
                      0 & 1 & 1 \\
                      0 & 0 & 1
                  \end{bmatrix}                                                        \\
                   & =\begin{bmatrix}\vecveps_1&\vecveps_2&\vecveps_3\end{bmatrix}\matb
              \end{align*}
              又因为\(\rank{\mata}=\rank{\matb}=3\)，由习题4.2可知\(\vecxi_1,\vecxi_2,\vecxi_3\)和\(\veceta_1,\veceta_2,\veceta_3\)均线性无关，且任意\(V\)中的向量均可由上述两个向量组线性表出。所以\(\vecxi_1,\vecxi_2,\vecxi_3\)和\(\veceta_1,\veceta_2,\veceta_3\)都是\(V\)的一组基；
              }
        \item {
              由(1)可得
              \begin{align*}
                  \begin{bmatrix}\veceta_1&\veceta_2&\veceta_3\end{bmatrix}
                   & =
                  \begin{bmatrix}\vecveps_1&\vecveps_2&\vecveps_3\end{bmatrix}\matb          \\
                   & =\begin{bmatrix}\vecxi_1&\vecxi_2&\vecxi_3\end{bmatrix}\inv{\mata}\matb
              \end{align*}
              因此由基\(\vecxi_1,\vecxi_2,\vecxi_3\)到基\(\veceta_1,\veceta_2,\veceta_3\)的过渡矩阵为
              \begin{equation*}
                  \inv{\mata}\matb=
                  \begin{bmatrix}
                      2  & 2  & 1 \\
                      2  & 3  & 1 \\
                      -1 & -1 & 0
                  \end{bmatrix}
              \end{equation*}
              }
    \end{enumerate}
\end{frame}

\subsection*{4.7}
\begin{frame}[allowframebreaks]
    \frametitle{题面}
    在线性空间\(\rea^{2\times2}\)中，已知
    \begin{equation*}
        \vecal_1=\begin{bmatrix}1&0\\0&0\end{bmatrix},
        \vecal_2=\begin{bmatrix}0&1\\0&0\end{bmatrix},
        \vecal_3=\begin{bmatrix}0&0\\1&0\end{bmatrix},
        \vecal_4=\begin{bmatrix}0&0\\0&1\end{bmatrix}
    \end{equation*}
    为其一组基，若\(\rea^{2\times2}\)的另一组基为\(\vecbeta_1,\vecbeta_2,\vecbeta_3,\vecbeta_4\)，由\(\vecal_1,\vecal_2,\vecal_3,\vecal_4\)到\(\vecbeta_1,\vecbeta_2,\vecbeta_3,\vecbeta_4\)的过渡矩阵为
    \begin{equation*}
        \mata=
        \begin{bmatrix}
            0 & 1 & 1 & 1 \\
            1 & 0 & 1 & 1 \\
            1 & 1 & 0 & 1 \\
            1 & 1 & 1 & 0
        \end{bmatrix}
    \end{equation*}
    求：
    \begin{enumerate}
        \item 基\(\vecbeta_1,\vecbeta_2,\vecbeta_3,\vecbeta_4\)；
        \item 矩阵\begin{equation*}\begin{bmatrix}0&1\\2&-3\end{bmatrix}\end{equation*}在基\(\vecbeta_1,\vecbeta_2,\vecbeta_3,\vecbeta_4\)下的坐标。
    \end{enumerate}
\end{frame}

\begin{frame}
    \frametitle{解答}
    \begin{enumerate}
        \item {
              由题设可知
              \begin{align*}
                  \vecbeta_1 & =\vecal_2+\vecal_3+\vecal_4=\begin{bmatrix}0&1\\1&1\end{bmatrix},
                  \vecbeta_2=\vecal_1+\vecal_3+\vecal_4=\begin{bmatrix}1&0\\1&1\end{bmatrix}     \\
                  \vecbeta_3 & =\vecal_1+\vecal_2+\vecal_4=\begin{bmatrix}1&1\\0&1\end{bmatrix},
                  \vecbeta_4=\vecal_1+\vecal_2+\vecal_3=\begin{bmatrix}1&1\\1&0\end{bmatrix}
              \end{align*}
              }\pause
        \item {
              容易得到该矩阵在基\(\vecal_1,\vecal_2,\vecal_3,\vecal_4\)下的坐标为\(\myvec{0,1,2,-3}\)。因此该矩阵在基\(\vecbeta_1,\vecbeta_2,\vecbeta_3,\vecbeta_4\)下的坐标为
              \begin{equation*}
                  \inv{\mata}\begin{bmatrix}0&1&2&-3\end{bmatrix}^\top=\begin{bmatrix}0&-1&-2&3\end{bmatrix}^\top
              \end{equation*}
              }
    \end{enumerate}
\end{frame}

\subsection*{4.8}
\begin{frame}
    \frametitle{题面}
    已知实数域\(\rea\)上的所有\(2\)阶矩阵，对于矩阵的加法和数乘，构成\(\rea\)上的四维线性空间，记作\(V=\rea^{2\times2}\)。
    \begin{enumerate}
        \item {
              分别证明
              \begin{equation*}
                  \vecal_1=\begin{bmatrix}1&0\\0&0\end{bmatrix},
                  \vecal_2=\begin{bmatrix}1&1\\0&0\end{bmatrix},
                  \vecal_3=\begin{bmatrix}1&1\\1&0\end{bmatrix},
                  \vecal_4=\begin{bmatrix}1&1\\1&1\end{bmatrix}
              \end{equation*}
              与
              \begin{equation*}
                  \vecbeta_1=\begin{bmatrix}-1&1\\1&1\end{bmatrix},
                  \vecbeta_2=\begin{bmatrix}1&-1\\1&1\end{bmatrix},
                  \vecbeta_3=\begin{bmatrix}1&1\\-1&1\end{bmatrix},
                  \vecbeta_4=\begin{bmatrix}1&1\\1&-1\end{bmatrix}
              \end{equation*}
              均为\(\rea^{2\times2}\)的基；}
        \item 对于\(\rea^{2\times2}\)中的任意元素\(\vecal\)，求\(\vecal\)在基\(\vecal_1,\vecal_2,\vecal_3,\vecal_4\)下的坐标；
        \item 求由基\(\vecal_1,\vecal_2,\vecal_3,\vecal_4\)到\(\vecbeta_1,\vecbeta_2,\vecbeta_3,\vecbeta_4\)的过渡矩阵。
    \end{enumerate}
\end{frame}

\begin{frame}[allowframebreaks]
    \frametitle{解答}
    \begin{enumerate}
        \item {
              由习题4.7可知，
              \begin{equation*}
                  \vecveps_1=\begin{bmatrix}1&0\\0&0\end{bmatrix},
                  \vecveps_2=\begin{bmatrix}0&1\\0&0\end{bmatrix},
                  \vecveps_3=\begin{bmatrix}0&0\\1&0\end{bmatrix},
                  \vecveps_4=\begin{bmatrix}0&0\\0&1\end{bmatrix}
              \end{equation*}
              为\(\rea^{2\times2}\)的一组基。又因为
              \begin{align*}
                  \begin{bmatrix}\vecal_1&\vecal_2&\vecal_3&\vecal_4\end{bmatrix}
                   & =
                  \begin{bmatrix}\vecveps_1&\vecveps_2&\vecveps_3&\vecveps_4\end{bmatrix}
                  \begin{bmatrix}
                      1 & 1 & 1 & 1 \\
                      0 & 1 & 1 & 1 \\
                      0 & 0 & 1 & 1 \\
                      0 & 0 & 0 & 1
                  \end{bmatrix} \\
                   & =
                  \begin{bmatrix}\vecveps_1&\vecveps_2&\vecveps_3&\vecveps_4\end{bmatrix}\mata
              \end{align*}
              \begin{align*}
                  \begin{bmatrix}\vecbeta_1&\vecbeta_2&\vecbeta_3&\vecbeta_4\end{bmatrix}
                   & =
                  \begin{bmatrix}\vecveps_1&\vecveps_2&\vecveps_3&\vecveps_4\end{bmatrix}
                  \begin{bmatrix}
                      -1 & 1  & 1  & 1  \\
                      1  & -1 & 1  & 1  \\
                      1  & 1  & -1 & 1  \\
                      1  & 1  & 1  & -1
                  \end{bmatrix}                                                                \\
                   & =\begin{bmatrix}\vecveps_1&\vecveps_2&\vecveps_3&\vecveps_4\end{bmatrix}\matb
              \end{align*}
              其中\(\rank{\mata}=\rank{\matb}=4\)。所以\(\vecal_1,\vecal_2,\vecal_3,\vecal_4\)和\(\vecbeta_1,\vecbeta_2,\vecbeta_3,\vecbeta_4\)线性无关，且\(\rea^{2\times2}\)中的任意元素可由这两个向量组线性表出。
              由基的定义可知上述两个向量组均为\(\rea^{2\times2}\)的基；
              }
        \item {
              设\(\vecal=\begin{bmatrix}\alpha_{11}&\alpha_{12}\\\alpha_{21}&\alpha_{22}\end{bmatrix}\)，则\(\vecal\)在基\(\vecveps_1,\vecveps_2,\vecveps_3,\vecveps_4\)下的坐标为\(\myvec{\alpha_{11},\alpha_{12},\alpha_{21},\alpha_{22}}\)。由(1)可得\(\vecal\)在基\(\vecal_1,\vecal_2,\vecal_3,\vecal_4\)下的坐标为
              \begin{equation*}
                  \inv{\mata}\begin{bmatrix}\alpha_{11}\\\alpha_{12}\\\alpha_{21}\\\alpha_{22}\end{bmatrix}=
                  \begin{bmatrix}\alpha_{11}-\alpha_{12}\\\alpha_{12}-\alpha_{21}\\\alpha_{21}-\alpha_{22}\\\alpha_{22}\end{bmatrix}
              \end{equation*}
              }
        \item {
              由(1)可得
              \begin{align*}
                  \begin{bmatrix}\vecbeta_1&\vecbeta_2&\vecbeta_3&\vecbeta_4\end{bmatrix}
                   & =\begin{bmatrix}\vecveps_1&\vecveps_2&\vecveps_3&\vecveps_4\end{bmatrix}\matb    \\
                   & =\begin{bmatrix}\vecal_1&\vecal_2&\vecal_3&\vecal_4\end{bmatrix}\inv{\mata}\matb
              \end{align*}
              所以由基\(\vecal_1,\vecal_2,\vecal_3,\vecal_4\)到\(\vecbeta_1,\vecbeta_2,\vecbeta_3,\vecbeta_4\)的过渡矩阵为
              \begin{equation*}
                  \inv{\mata}\matb=
                  \begin{bmatrix}
                      -2 & 2  & 0  & 0  \\
                      0  & -2 & 2  & 0  \\
                      0  & 0  & -2 & 2  \\
                      1  & 1  & 1  & -1
                  \end{bmatrix}
              \end{equation*}
              }
    \end{enumerate}
\end{frame}

\subsection*{4.9}
\begin{frame}
    \frametitle{题面}
    求下列两个齐次线性方程组的解空间的基和维数：
    \begin{enumerate}
        \item \(x_1+x_2+\cdots+x_n=0\)；
        \item \begin{equation*}
                  \begin{cases}
                      2x_1-4x_2+5x_3+3x_4=0 \\
                      3x_1-6x_2+4x_3+2x_4=0 \\
                      4x_1-8x_2+17x_3+11x_4=0
                  \end{cases}
              \end{equation*}
    \end{enumerate}
\end{frame}

\begin{frame}[allowframebreaks]
    \frametitle{解答}
    \begin{enumerate}
        \item {
              方程组对应的系数矩阵的最简阶梯型矩阵为\(\begin{bmatrix}1&1&\cdots&1\end{bmatrix}\)，因此解空间的基为
              \begin{equation*}
                  \begin{bmatrix}-1\\1\\0\\0\\\vdots\\0\\0\end{bmatrix},
                  \begin{bmatrix}-1\\0\\1\\0\\\vdots\\0\\0\end{bmatrix},\cdots
                  \begin{bmatrix}-1\\0\\0\\0\\\vdots\\0\\1\end{bmatrix}
              \end{equation*}
              维数为\(n-1\)；
              }
        \item {
              方程组对应的系数矩阵的最简阶梯型矩阵为
              \begin{equation*}
                  \begin{bmatrix}
                      1 & -2 & 0 & -\frac{2}{7} \\
                      0 & 0  & 1 & \frac{5}{7}  \\
                      0 & 0  & 0 & 0
                  \end{bmatrix}
              \end{equation*}
              因此解空间的基为\(\myvec{2,1,0,0}\)，\(\myvec{2,0,-5,7}\)，维数为\(2\)。
              }
    \end{enumerate}
\end{frame}

\subsection*{4.10}
\begin{frame}
    \frametitle{题面}
    设\(\enums{\vecal}{n}\)是\(n\)维线性空间\(V\)的一组基，又\(V\)中向量\(\vecal_{n+1}\)在这组基下坐标\(\mypar{\enums{x}{n}}\)全都不为零。

    证明\(\enums{\vecal}{n},\vecal_{n+1}\)中任意\(n\)个向量必构成\(V\)的一组基，并求\(\vecal_1\)在基\(\vecal_2,\dots,\vecal_n,\vecal_{n+1}\)下的坐标。
\end{frame}

\begin{frame}[allowframebreaks]
    \frametitle{解答}
    因为\(\enums{\vecal}{n}\)是\(V\)的一组基，下面证明对\(\forall i=1,2,\dots,n\)，
    \begin{equation*}
        \mata_i=\setof{\vecal_1,\dots,\vecal_{i-1},\vecal_{i+1},\dots,\vecal_{n+1}}
    \end{equation*}
    都是\(V\)的一组基。

    由题设可知
    \begin{align*}
        \begin{bmatrix}\vecal_1\\\vdots\\\vecal_{i+1}\\\vdots\\\vecal_{n+1}\end{bmatrix}^\top
         & =
        \begin{bmatrix}\vecal_1\\\vecal_2\\\vdots\\\vecal_n\end{bmatrix}^\top
        \begin{bmatrix}
            1      & \cdots & 0      & 0      & \cdots & 0      & x_1     \\
            \vdots & \ddots & \vdots & \vdots & \cdots & \vdots & \vdots  \\
            0      & \cdots & 1      & 0      & \cdots & 0      & x_{i-1} \\
            0      & \cdots & 0      & 0      & \cdots & 0      & x_i     \\
            0      & \cdots & 0      & 1      & \cdots & 0      & x_{i+1} \\
            \vdots &        & \vdots & \vdots & \ddots & \vdots & \vdots  \\
            0      & \cdots & 0      & 0      & \cdots & 1      & x_n
        \end{bmatrix} \\
         & =
        \begin{bmatrix}\vecal_1\\\vecal_2\\\vdots\\\vecal_n\end{bmatrix}^\top\matm_i
    \end{align*}
    容易发现\(\rank{\matm_i}=n\)，所以\(\mata_i\)线性无关，且\(V\)中任意向量可由\(\mata_i\)线性表出。

    因此\(\enums{\vecal}{n},\vecal_{n+1}\)中任意\(n\)个向量必构成\(V\)的一组基。

    同时，由等式\(\vecal_{n+1}=x_1\vecal_1+x_2\vecal_2+\cdots+x_n\vecal_n\)可得
    \begin{equation*}
        \vecal_1=-\frac{x_2}{x_1}\vecal_2-\cdots-\frac{x_n}{x_1}\vecal_n+\frac{1}{x_1}\vecal_{n+1}
    \end{equation*}
    所以\(\vecal_1\)在在基\(\vecal_2,\dots,\vecal_n,\vecal_{n+1}\)下的坐标为\(\myvec{-\frac{x_2}{x_1},\cdots,-\frac{x_n}{x_1},\frac{1}{x_1}}\)。
\end{frame}

\subsection*{4.12}
\begin{frame}
    \frametitle{题面}
    设\(V_1,V_2\)是\(\rea^n\)的两个非平凡子空间，证明：在\(\rea^n\)中存在向量\(\vecal\)，使\(\vecal\notin V_1\)，且\(\vecal\notin V_2\)，并在\(\rea^3\)中举例说明此结论。
\end{frame}

\begin{frame}
    \frametitle{解答}
    \begin{enumerate}
        \item {
              若\(V_1,V_2\)其中一个是另一个的子空间，则不妨设\(V_1\subseteq V_2\)。

              因为\(V_2\)非平凡，存在\(\vecal\in\rea^n\)满足\(\vecal\notin V_2\)，从而\(\vecal\notin V_1\)；
              }\pause
        \item {
              若\(V_1,V_2\)中任意空间都不是另一个的子空间，则存在\(\vecal_1\in V_1\cap\overline{V_2}\)，\(\vecal_2\in V_2\cap\overline{V_1}\)。此时构造\(\vecal=\vecal_1+\vecal_2\)：

              若\(\vecal\in V_1\)，则有\(\vecal_2=\vecal-\vecal_1\in V_1\)，与假设矛盾。因此\(\vecal\notin V_1\)。

              同理可得\(\vecal\notin V_2\)，结论得证。

              举例：\(V_1=\myspan{\myvec{1,0,0},\myvec{0,1,0}}\)，\(V_2=\myspan{\myvec{0,1,0},\myvec{0,0,1}}\)，\(\vecal=\myvec{1,0,1}\)。
              }
    \end{enumerate}
\end{frame}

\subsection*{4.13}
\begin{frame}
    \frametitle{题面}
    设\(\vecal,\vecbeta,\vecgamma\in\rea^n\)，\(c_1,c_2,c_3\in\rea\)，且\(c_1c_3\neq0\)，证明：若\(c_1\vecal+c_2\vecbeta+c_3\vecgamma=\veczero\)，则\(\myspan{\vecal,\vecbeta}=\myspan{\vecbeta,\vecgamma}\)。
\end{frame}

\begin{frame}
    \frametitle{解答}
    因为\(c_1\vecal+c_2\vecbeta+c_3\vecgamma=\veczero\)，且\(c_1c_3\neq0\)，所以有\(\vecal=-\frac{c_2}{c_1}\vecbeta-\frac{c_3}{c_1}\vecgamma\)。

    因此对\(\forall\veceta\in\myspan{\vecal,\vecbeta}\)，不妨设\(\veceta=\lambda_1\vecal+\lambda_2\vecbeta\)，则有
    \begin{align*}
        \veceta & =\lambda_1\mypar{-\frac{c_2}{c_1}\vecbeta-\frac{c_3}{c_1}\vecgamma}+\lambda_2\vecbeta \\
                & =\mypar{\lambda_2-\frac{c_2\lambda_1}{c_1}}\vecbeta-\frac{c_3\lambda_1}{c_1}\vecgamma
    \end{align*}
    所以\(\veceta\in\myspan{\vecbeta,\vecgamma}\)，即\(\myspan{\vecal,\vecbeta}\subseteq\myspan{\vecbeta,\vecgamma}\)。

    同理可得\(\myspan{\vecbeta,\vecgamma}\subseteq\myspan{\vecal,\vecbeta}\)。因此\(\myspan{\vecal,\vecbeta}=\myspan{\vecbeta,\vecgamma}\)。
\end{frame}

\subsection*{4.14}
\begin{frame}
    \frametitle{题面}
    若
    \begin{equation*}
        \mata=\begin{bmatrix}1&1&1&0\\2&1&0&1\end{bmatrix}
    \end{equation*}
    求\(\mata\)的零空间\(\nullsp{\mata}\)的一组基。
\end{frame}

\begin{frame}
    \frametitle{解答}
    对\(A\)使用初等变换，将其化为最简阶梯型矩阵：
    \begin{equation*}
        \begin{bmatrix}
            1 & 0 & -1 & 1  \\
            0 & 1 & 2  & -1
        \end{bmatrix}
    \end{equation*}
    因此得到\(\nullsp{\mata}\)的一组基为\(\myvec{1,-2,1,0}\)，\(\myvec{-1,1,0,1}\)。
\end{frame}

\subsection*{4.16}
\begin{frame}
    \frametitle{题面}
    已知向量组\(\vecal_1=\myvec{2,0,1,3,-1}\)，\(\vecal_2=\myvec{0,-2,1,5,-3}\)，\(\vecbeta_1=\myvec{1,1,0,-1,1}\)，\(\vecbeta_2=\myvec{1,-3,2,0,5}\)，且\(W_1=\myspan{\vecal_1,\vecal_2}\)，\(W_2=\myspan{\vecbeta_1,\vecbeta_2}\)。试求\(W_1\cap W_2\)与\(W_1+W_2\)的维数以及各自的一组基。
\end{frame}

\begin{frame}
    \frametitle{解答}
    因为\(W_1+W_2=\myspan{\vecal_1,\vecal_2,\vecbeta_1,\vecbeta_2}\)，因此考虑对\(\mat{\vecal_1,\vecal_2,\vecbeta_1,\vecbeta_2}\)作初等行变换，化为最简阶梯型矩阵：
    \begin{equation*}
        \begin{bmatrix}
            2  & 0  & 1  & 1  \\
            0  & -2 & 1  & -3 \\
            1  & 1  & 0  & 2  \\
            3  & 5  & -1 & 0  \\
            -1 & -3 & 1  & 5
        \end{bmatrix}\Rightarrow
        \begin{bmatrix}
            1 & 0 & \frac{1}{2}  & 0 \\
            0 & 1 & -\frac{1}{2} & 0 \\
            0 & 0 & 0            & 1 \\
            0 & 0 & 0            & 0 \\
            0 & 0 & 0            & 0
        \end{bmatrix}
    \end{equation*}
    所以\(W_1+W_2\)的维数为\(3\)，一组基为\(\vecal_1,\vecal_2,\vecbeta_2\)。

    由维数公式可知\(W_1\cap W_2\)的维数为\(1\)，又因为\(\vecbeta_1=\frac{1}{2}\vecal_1-\frac{1}{2}\vecal_2\in W_1\)，所以\(\vecbeta_1\in W_1\cap W_2\)为其一组基。
\end{frame}

\subsection*{4.17}
\begin{frame}
    \frametitle{题面}
    设\(V\)是内积空间，试证明：
    \begin{enumerate}
        \item 对\(\forall\vecal\in V\)，都有\(\indot{\veczero}{\vecal}=0\)；
        \item 若对\(\forall\vecal\in V\)，恒有\(\indot{\vecal}{\vecbeta}=0\)，则\(\vecbeta=\veczero\)；
        \item 已知\(\vecal,\vecbeta\in V\)，如对\(\forall\vecgamma\in V\)，都有\(\indot{\vecal}{\vecgamma}=\indot{\vecbeta}{\vecgamma}\)。试证明\(\vecal=\vecbeta\)。
    \end{enumerate}
\end{frame}

\begin{frame}
    \frametitle{解答}
    \begin{enumerate}
        \item {
              对\(\forall\vecbeta\in V\)，有
              \begin{equation*}
                  \indot{\veczero}{\vecal}=\indot{\vecbeta+\mypar{-\vecbeta}}{\vecal}=\indot{\vecbeta}{\vecal}-\indot{\vecbeta}{\vecal}=0
              \end{equation*}
              }\pause
        \item 令\(\vecal=\vecbeta\)，得到\(\indot{\vecbeta}{\vecbeta}=0\)，由内积定义知\(\vecbeta=\veczero\)；\pause
        \item 由\(\indot{\vecal}{\vecgamma}=\indot{\vecbeta}{\vecgamma}\)可得\(\indot{\vecal-\vecbeta}{\vecgamma}=0\)。由\(\vecgamma\)的任意性及(2)结论可知\(\vecal-\vecbeta=\veczero\)，即\(\vecal=\vecbeta\)。
    \end{enumerate}
\end{frame}

\subsection*{4.20}
\begin{frame}
    \frametitle{题面}
    设\(\enums{\vecal}{n}\)是内积空间的\(n\)个向量，行列式
    \begin{equation*}
        \gram{\enums{\vecal}{n}}=
        \begin{vmatrix}
            \indot{\vecal_1}{\vecal_1} & \indot{\vecal_1}{\vecal_2} & \cdots & \indot{\vecal_1}{\vecal_n} \\
            \indot{\vecal_2}{\vecal_1} & \indot{\vecal_2}{\vecal_2} & \cdots & \indot{\vecal_2}{\vecal_n} \\
            \vdots                     & \vdots                     &        & \vdots                     \\
            \indot{\vecal_n}{\vecal_1} & \indot{\vecal_n}{\vecal_2} & \cdots & \indot{\vecal_n}{\vecal_n}
        \end{vmatrix}
    \end{equation*}
    称为\(\enums{\vecal}{n}\)的Gram行列式。证明\(\gram{\enums{\vecal}{n}}=0\)的充分必要条件是\(\enums{\vecal}{n}\)线性相关。
\end{frame}

\begin{frame}[allowframebreaks]
    \frametitle{解答}
    \begin{itemize}
        \item {
              首先证明充分性。

              若\(\enums{\vecal}{n}\)线性相关，则有\(\sum_{i=1}^nk_i\vecal_i=\veczero\)，其中\(k_i\)不全为零。
              因此对\(\forall j=1,2,\dots,n\)，有
              \begin{equation*}
                  \indot{\vecal_j}{\sum_{i=1}^nk_i\vecal_i}=\sum_{i=1}^nk_i\indot{\vecal_j}{\vecal_i}=0
              \end{equation*}
              所以\(\gram{\enums{\vecal}{n}}\)对应矩阵的列线性相关，进一步得到
              \begin{equation*}
                  \gram{\enums{\vecal}{n}}=0
              \end{equation*}
              }
        \item {
              随后证明必要性。
              若\(\gram{\enums{\vecal}{n}}=0\)，则其对应矩阵的列线性相关。因此对\(\forall j=1,2,\dots,n\)，有\(\sum_{i=1}^nk_i\indot{\vecal_j}{\vecal_i}=0\)，其中\(k_i\)不全为零，从而进一步得到\(\indot{\vecal_j}{\sum_{i=1}^nk_i\vecal_i}=0\)。由内积性质得到
              \begin{equation*}
                  \indot{\sum_{i=1}^nk_i\vecal_i}{\sum_{i=1}^nk_i\vecal_i}=0
              \end{equation*}
              因此\(\sum_{i=1}^nk_i\vecal_i=\veczero\)，\(\enums{\vecal}{n}\)线性相关。
              }
    \end{itemize}
\end{frame}

\subsection*{4.21}
\begin{frame}
    \frametitle{题面}
    设\(\enums{\vecal}{r}\)线性无关，\(\vecbeta_1,\vecbeta_2\)也线性无关，且后者与前者中的每个向量都正交，证明：\(\vecbeta_1,\vecbeta_2,\enums{\vecal}{r}\)也线性相关。
\end{frame}

\begin{frame}[allowframebreaks]
    \frametitle{解答}
    反设结论不成立，则存在一组不全为零的数\(l_1,l_2,k_1,k_2,\dots,k_r\)满足
    \begin{equation*}
        l_1\vecbeta_1+l_2\vecbeta_2+\sum_{i=1}^rk_i\vecal_i=\veczero
    \end{equation*}
    此时\(l_1,l_2\)不全为零，否则\(\enums{\vecal}{r}\)线性相关，与题设矛盾。

    又因为\(\vecbeta_1,\vecbeta_2\)和\(\enums{\vecal}{r}\)中的每个向量正交，因此有
    \begin{align*}
          & \indot{l_1\vecbeta_1+l_2\vecbeta_2+\sum_{i=1}^rk_i\vecal_i}{l_1\vecbeta_1+l_2\vecbeta_2}   \\
        = & \indot{l_1\vecbeta_1+l_2\vecbeta_2}{l_1\vecbeta_1+l_2\vecbeta_2}
        +l_1\sum_{i=1}^r\indot{k_i\vecal_i}{\vecbeta_1}+l_2\sum_{i=1}^r\indot{k_i\vecal_i}{\vecbeta_1} \\
        = & \indot{l_1\vecbeta_1+l_2\vecbeta_2}{l_1\vecbeta_1+l_2\vecbeta_2}=0
    \end{align*}
    由内积定义知\(l_1\vecbeta_1+l_2\vecbeta_2=\veczero\)，即\(\vecbeta_1,\vecbeta_2\)线性相关。这与题设矛盾，因此结论成立。
\end{frame}

\subsection*{4.22}
\begin{frame}
    \frametitle{题面}

    试把\(\myvec{1,0,1,0}\)，\(\myvec{0,1,0,2}\)扩充成为\(\rea^4\)的一组标准正交基。

\end{frame}

\begin{frame}
    \frametitle{解答}
    设\(\vecal_1=\myvec{1,0,1,0}\)，\(\vecal_2=\myvec{0,1,0,2}\)。容易发现\(\vecal_1,\vecal_2\)正交。

    设\(\rea^4\)的剩余正交基为\(\vecx\)，则有\(\begin{bmatrix}\vecal_1^\top\\\vecal_2^\top\end{bmatrix}\vecx=\veczero\)。

    对该齐次线性方程组求解，得到基础解系：\(\vecal_3=\myvec{-1,0,1,0}\)，\(\vecal_4=\myvec{0,-2,0,1}\)。

    将上述正交基标准化，得到\(\rea^4\)的一组标准正交基：

    \begin{equation*}
        \begin{bmatrix}
            \vecq_1 & \vecq_2 & \vecq_3 & \vecq_4
        \end{bmatrix}=
        \begin{bmatrix}
            \frac{\sqrt{2}}{2} & 0                   & -\frac{\sqrt{2}}{2} & 0                    \\
            0                  & \frac{\sqrt{5}}{5}  & 0                   & -\frac{2\sqrt{5}}{5} \\
            \frac{\sqrt{2}}{2} & 0                   & \frac{\sqrt{2}}{2}  & 0                    \\
            0                  & \frac{2\sqrt{5}}{5} & 0                   & \frac{\sqrt{5}}{5}
        \end{bmatrix}
    \end{equation*}
\end{frame}

\subsection*{4.23}
\begin{frame}
    \frametitle{题面}
    已知\(f_1=1\)，\(f_2=x-1\)，\(f_3=\mypar{x-1}^2\)为三维内积空间\(C\mat{0,1}\)的一组基，内积定义为
    \begin{equation*}
        \indot{f}{g}=\int_0^1f(x)g(x)\mathrm{d}x
    \end{equation*}
    利用Schmidt正交化方法，求与\(f_1,f_2,f_3\)等价的一组标准正交基（两个向量组可相互线性表出称为等价）。
\end{frame}

\begin{frame}
    \frametitle{解答}
    \begin{align*}
        \vecg_1 & =\frac{\vecf_1}{\Abs{\vecf_1}}=\frac{\vecf_1}{\indot{\vecf_1}{\vecf_1}}=1                                                                                      \\
        \vecg_2 & =\frac{\vecf_2-\indot{\vecg_1}{\vecf_2}\vecg_1}{\Abs{\vecf_2-\indot{\vecg_1}{\vecf_2}\vecg_1}}                                                                 \\
                & =\frac{x-\frac{1}{2}}{\sqrt{\frac{1}{12}}}=\sqrt3\mypar{2x-1}                                                                                                  \\
        \vecg_3 & =\frac{\vecf_3-\indot{\vecg_1}{\vecf_3}\vecg_1-\indot{\vecg_2}{\vecf_3}\vecg_2}{\Abs{\vecf_3-\indot{\vecg_1}{\vecf_3}\vecg_1-\indot{\vecg_2}{\vecf_3}\vecg_2}} \\
                & =\frac{x^2-x+\frac{1}{6}}{\sqrt{\frac{1}{180}}}=\sqrt5\mypar{6x^2-6x+1}
    \end{align*}
\end{frame}

\subsection*{4.24}
\begin{frame}
    \frametitle{题面}
    设\(\enums{\vecal}{n}\)与\(\enums{\vecbeta}{n}\)为\(n\)维欧氏空间的两组基，其中\(\enums{\vecal}{n}\)为标准正交基，又若由\(\enums{\vecal}{n}\)到\(\enums{\vecbeta}{n}\)的过渡矩阵为\(\matm\)，证明：\(\enums{\vecbeta}{n}\)为\(V\)的标准正交基的充分必要条件是\(\matm\)为正交矩阵。
\end{frame}

\begin{frame}
    \frametitle{解答}
    设\(\mata=\mat{\enums{\vecal}{n}}\)，\(\matb=\mat{\enums{\vecbeta}{n}}\)。

    因为\(\enums{\vecal}{n}\)为标准正交基，所以\(\mata\)满足\(\mata^\top\mata=\mati\)。

    又由过渡矩阵的定义知\(\matb=\mata\matm\)，因此有

    \begin{align*}
        \matb^\top\matb & =\mypar{\mata\matm}^\top\mata\matm      \\
                        & =\matm^\top\mypar{\mata^\top\mata}\matm \\
                        & =\matm^\top\matm
    \end{align*}
    因此\(\matm\)为正交矩阵等价于\(\matb\)为正交矩阵。又因为\(\matb\)为正交矩阵等价于\(\enums{\vecbeta}{n}\)为\(V\)的标准正交基，所以\(\enums{\vecbeta}{n}\)为\(V\)的标准正交基的充分必要条件是\(\matm\)为正交矩阵。
\end{frame}

\subsection*{4.26}
\begin{frame}
    \frametitle{题面}
    设
    \begin{equation*}
        \mata=
        \begin{bmatrix}
            1 & 2  & -1 \\
            2 & 0  & 1  \\
            2 & -4 & 2  \\
            4 & 0  & 0
        \end{bmatrix},
        \vecb=\begin{bmatrix}-1\\1\\1\\-2\end{bmatrix}
    \end{equation*}
    问线性方程组\(\mata\vecx=\vecb\)是否有解？如果无解，求其最小二乘解。
\end{frame}

\begin{frame}
    \frametitle{解答}
    分别对\(\mata\)和\(\mat{\mata,\vecb}\)作初等行变换，可得\(\rank{\mata}=3\)，\(\rank{\mat{\mata,\vecb}}=4\)。所以线性方程组\(\mata\vecx=\vecb\)无解。

    下面考虑求方程的最小二乘解：令
    \begin{equation*}
        \mata'=\mata^\top\mata=
        \begin{bmatrix}
            25 & -6  & 5   \\
            -6 & 20  & -10 \\
            5  & -10 & 6
        \end{bmatrix}
    \end{equation*}
    \begin{equation*}
        \vecb'=\mata^\top\vecb=\begin{bmatrix}-5&6&4\end{bmatrix}^\top
    \end{equation*}
    解线性方程组\(\mata'\hat{\vecx}=\vecb'\)，得到最小二乘解为
    \begin{equation*}
        \hat{\vecx}=
        \begin{bmatrix}
            -\frac{11}{24} & \frac{25}{48} & \frac{23}{12}
        \end{bmatrix}^\top
    \end{equation*}
\end{frame}

\subsection*{4.27}
\begin{frame}
    \frametitle{题面}
    求向量\(\vecb\)在矩阵\(\mata\)的列空间中的正交投影。
    \begin{equation*}
        \mata=
        \begin{bmatrix}
            1 & -2 \\
            1 & 0  \\
            1 & 1  \\
            1 & 3
        \end{bmatrix},
        \vecb=\begin{bmatrix}-4\\-3\\3\\0\end{bmatrix}
    \end{equation*}
\end{frame}

\begin{frame}
    \frametitle{解答}
    分别对\(\mata\)和\(\mat{\mata,\vecb}\)作初等行变换，可得\(\rank{\mata}=2\)，\(\rank{\mat{\mata,\vecb}}=3\)。所以\(\vecb\notin\colsp{\mata}\)。

    由例4.4.11可知\(\vecb\)在\(\colsp{\mata}\)上的正交投影为
    \begin{equation*}
        \hat{\vecb}=\mata\pinv{\mata^\top\mata}\mata^\top\vecb=
        \begin{bmatrix}
            -\frac{7}{2} & -\frac{3}{2} & -\frac{1}{2} & \frac{3}{2}
        \end{bmatrix}^\top
    \end{equation*}
\end{frame}

\subsection*{4.28}
\begin{frame}
    \frametitle{题面}
    设齐次线性方程组
    \begin{equation*}
        \begin{cases}
            2x_1+x_2-x_3+x_4-3x_5=0 \\
            x_1+x_2-x_3+x_5=0
        \end{cases}
    \end{equation*}
    \begin{enumerate}
        \item 求解空间的标准正交基；
        \item 求解空间的正交补空间的标准正交基。
    \end{enumerate}
\end{frame}

\begin{frame}[allowframebreaks]
    \frametitle{解答}
    记该方程组为\(\mata\vecx=\veczero\)，其中\(\mata\)为系数矩阵。
    \begin{enumerate}
        \item {
              将\(\mata\)通过初等行变换化为最简阶梯型矩阵：
              \begin{equation*}
                  \begin{bmatrix}
                      1 & 0 & 0  & 1  & -4 \\
                      0 & 1 & -1 & -1 & 5
                  \end{bmatrix}
              \end{equation*}
              从而求得\(\nullsp{\mata}\)的一组基为
              \begin{equation*}
                  \vecal_1=\begin{bmatrix}0\\1\\1\\0\\0\end{bmatrix},
                  \vecal_2=\begin{bmatrix}-1\\1\\0\\1\\0\end{bmatrix},
                  \vecal_3=\begin{bmatrix}4\\-5\\0\\0\\1\end{bmatrix}
              \end{equation*}
              将\(\vecal_1,\vecal_2,\vecal_3\) Schmidt正交化：
              \begin{align*}
                  \vecq_1 & =\frac{\vecal_1}{\Abs{\vecal_1}}=\frac{1}{\sqrt2}\begin{bmatrix}0&1&1&0&0\end{bmatrix}^\top                                                                                                                      \\
                  \vecq_2 & =\frac{\vecal_2-\vecq_1\vecq_1^\top\vecal_2}{\Abs{\vecal_2-\vecq_1\vecq_1^\top\vecal_2}}=\frac{1}{\sqrt{10}}\begin{bmatrix}-2&1&-1&2&0\end{bmatrix}^\top                                                         \\
                  \vecq_3 & =\frac{\vecal_3-\vecq_1\vecq_1^\top\vecal_3-\vecq_2\vecq_2^\top\vecal_3}{\Abs{\vecal_3-\vecq_1\vecq_1^\top\vecal_3-\vecq_2\vecq_2^\top\vecal_3}}=\frac{1}{3\sqrt{35}}\begin{bmatrix}4&-5&0&0&1\end{bmatrix}^\top
              \end{align*}
              \(\vecq_1,\vecq_2,\vecq_3\)构成\(\nullsp{\mata}\)的一组标准正交基。
              }
        \item {
              由推论4.4.1可知\(\nullsp{\mata}\)的正交补空间为\(\colsp{\mata^\top}\)，即\(\mata\)的行空间。

              由(1)可知\(\dim\mypar{\colsp{\mata^\top}}=2\)，因此\(\colsp{\mata^\top}\)的一组基为
              \begin{equation*}
                  \vecbeta_1=\begin{bmatrix}2&1&-1&1&-3\end{bmatrix}^\top,
                  \vecbeta_2=\begin{bmatrix}1&1&-1&0&1\end{bmatrix}^\top
              \end{equation*}
              同(1)理，将\(\vecbeta_1,\vecbeta_2\) Schmidt正交化后得到\(\colsp{\mata^\top}\)的标准正交基为
              \begin{align*}
                  \vecp_1 & =\frac{1}{4}\begin{bmatrix}2&1&-1&1&-3\end{bmatrix}^\top              \\
                  \vecp_2 & =\frac{1}{12\sqrt{7}}\begin{bmatrix}14&15&-15&-1&19\end{bmatrix}^\top
              \end{align*}
              }
    \end{enumerate}
\end{frame}

\subsection*{4.29}
\begin{frame}
    \frametitle{题面}
    设\(\vecal_1,\vecal_2\)和\(\vecal\)是Euclid空间\(\rea^3\)中的向量，
    \begin{equation*}
        \vecal_1=\begin{bmatrix}1\\1\\2\end{bmatrix},
        \vecal_2=\begin{bmatrix}1\\0\\1\end{bmatrix},
        \vecal=\begin{bmatrix}3\\2\\1\end{bmatrix}
    \end{equation*}
    求\(\myspan{\vecal_1,\vecal_2}\)的正交补空间以及向量\(\vecal\)分别在\(\myspan{\vecal_1,\vecal_2}\)及其正交补空间中的正交投影。
\end{frame}

\begin{frame}
    \frametitle{解答}
    记\(\mata=\mat{\vecal_1,\vecal_2}\)。
    \begin{itemize}
        \item {
              容易发现\(\vecal_1,\vecal_2\)线性无关，因此它们构成\(\myspan{\vecal_1,\vecal_2}\)的一组基。

              由例4.4.11可知\(\vecal\)在\(\myspan{\vecal_1,\vecal_2}\)中的正交投影为
              \begin{equation*}
                  \mata\pinv{\mata^\top\mata}\mata^\top\vecal=\begin{bmatrix}\frac{5}{3}&\frac{2}{3}&\frac{7}{3}\end{bmatrix}^\top
              \end{equation*}
              }
        \item {
              由推论4.4.1可知\(\myspan{\vecal_1,\vecal_2}\)的正交补空间为\(\nullsp{\mata^\top}\)。

              将\(\mata^\top\)通过初等行变换化为最简阶梯型矩阵\(\begin{bmatrix}1 & 0 & 1 \\0 & 1 & 1\end{bmatrix}\)，从而求得\(\nullsp{\mata^\top}\)的基\(\vecbeta_1=\myvec{-1,-1,1}\)。

              所以\(\vecal\)在\(\nullsp{\mata^\top}\)中的正交投影为
              \begin{equation*}
                  \frac{\vecbeta\vecbeta^\top}{\vecbeta^\top\vecbeta}\vecal=\begin{bmatrix}\frac{4}{3}&\frac{4}{3}&-\frac{4}{3}\end{bmatrix}^\top
              \end{equation*}
              }
    \end{itemize}
\end{frame}

\subsection*{4.31}
\begin{frame}
    \frametitle{题面}
    已知在\(\rea^3\)中，线性变换\(T\)在基
    \begin{equation*}
        \vecveps_1=\begin{bmatrix}8\\-6\\7\end{bmatrix},
        \vecveps_2=\begin{bmatrix}-16\\7\\-13\end{bmatrix},
        \vecveps_3=\begin{bmatrix}9\\-3\\7\end{bmatrix}
    \end{equation*}
    下的表示矩阵
    \begin{equation*}
        \mata=
        \begin{bmatrix}
            1  & -18 & 15 \\
            -1 & -22 & 20 \\
            1  & -25 & 22
        \end{bmatrix}
    \end{equation*}
    求\(T\)在基
    \begin{equation*}
        \veceta_1=\begin{bmatrix}1\\-2\\1\end{bmatrix},
        \veceta_2=\begin{bmatrix}3\\-1\\2\end{bmatrix},
        \veceta_3=\begin{bmatrix}2\\1\\2\end{bmatrix}
    \end{equation*}
    下的表示矩阵。
\end{frame}

\begin{frame}
    \frametitle{解答}
    设从基\(\vecveps_1,\vecveps_2,\vecveps_3\)到基\(\veceta_1,\veceta_2,\veceta_3\)的过渡矩阵为\(\matm\)，则有
    \begin{equation*}
        \begin{bmatrix}\veceta_1&\veceta_2&\veceta_3\end{bmatrix}=
        \begin{bmatrix}\vecveps_1&\vecveps_2&\vecveps_3\end{bmatrix}\matm
    \end{equation*}
    \begin{equation*}
        \matm=\inv{\begin{bmatrix}\vecveps_1&\vecveps_2&\vecveps_3\end{bmatrix}}\begin{bmatrix}\veceta_1&\veceta_2&\veceta_3\end{bmatrix}=
        \begin{bmatrix}
            1 & 1 & -3 \\
            1 & 2 & -5 \\
            1 & 3 & -6
        \end{bmatrix}
    \end{equation*}
    由定理4.5.4知\(T\)在基\(\veceta_1,\veceta_2,\veceta_3\)下的表示矩阵为
    \begin{equation*}
        \inv{\matm}\mata\matm=
        \begin{bmatrix}
            1 & 2  & 2  \\
            3 & -1 & -2 \\
            2 & -3 & 1
        \end{bmatrix}
    \end{equation*}
\end{frame}

\subsection*{4.33}
\begin{frame}
    \frametitle{题面}
    给定\(\rea^3\)的两组基：
    \begin{equation*}
        \vecveps_1=\begin{bmatrix}1\\0\\1\end{bmatrix},
        \vecveps_2=\begin{bmatrix}2\\1\\0\end{bmatrix},
        \vecveps_3=\begin{bmatrix}1\\1\\1\end{bmatrix}
    \end{equation*}
    和
    \begin{equation*}
        \veceta_1=\begin{bmatrix}1\\2\\-1\end{bmatrix},
        \veceta_2=\begin{bmatrix}2\\2\\-1\end{bmatrix},
        \veceta_3=\begin{bmatrix}2\\-1\\-1\end{bmatrix}
    \end{equation*}
    定义线性变换：\(T\mypar{\vecveps_i}=\veceta_i\)，\(i\in\setof{1,2,3}\)。分别求\(T\)在基\(\vecveps_1,\vecveps_2,\vecveps_3\)与\(\veceta_1,\veceta_2,\veceta_3\)下的表示矩阵。
\end{frame}

\begin{frame}[allowframebreaks]
    \frametitle{解答}
    设从基\(\vecveps_1,\vecveps_2,\vecveps_3\)到基\(\veceta_1,\veceta_2,\veceta_3\)的过渡矩阵为\(\matm\)。则有
    \begin{equation*}
        \begin{bmatrix}\veceta_1&\veceta_2&\veceta_3\end{bmatrix}=
        \begin{bmatrix}\vecveps_1&\vecveps_2&\vecveps_3\end{bmatrix}\matm
    \end{equation*}
    \begin{equation*}
        \matm=\inv{\begin{bmatrix}\vecveps_1&\vecveps_2&\vecveps_3\end{bmatrix}}\begin{bmatrix}\veceta_1&\veceta_2&\veceta_3\end{bmatrix}=
        \begin{bmatrix}
            -2 & -\frac{3}{2} & \frac{3}{2}  \\
            1  & \frac{3}{2}  & \frac{3}{2}  \\
            1  & \frac{1}{2}  & -\frac{5}{2}
        \end{bmatrix}
    \end{equation*}
    设\(T\)在基\(\vecveps_1,\vecveps_2,\vecveps_3\)下的表示矩阵为\(\mata\)。由题设可知
    \begin{equation*}
        \begin{bmatrix}\veceta_1&\veceta_2&\veceta_3\end{bmatrix}=T\mypar{\vecveps_1,\vecveps_2,\vecveps_3}=\begin{bmatrix}\vecveps_1&\vecveps_2&\vecveps_3\end{bmatrix}\mata
    \end{equation*}
    因此有\(\mata=\matm\)。由定理4.5.4知由定理4.5.4知\(T\)在基\(\veceta_1,\veceta_2,\veceta_3\)下的表示矩阵为\(\inv{\matm}\mata\matm=\matm\)。

    综上，\(T\)在基\(\vecveps_1,\vecveps_2,\vecveps_3\)与\(\veceta_1,\veceta_2,\veceta_3\)下的表示矩阵均为
    \begin{equation*}
        \matm=
        \begin{bmatrix}
            -2 & -\frac{3}{2} & \frac{3}{2}  \\
            1  & \frac{3}{2}  & \frac{3}{2}  \\
            1  & \frac{1}{2}  & -\frac{5}{2}
        \end{bmatrix}
    \end{equation*}
\end{frame}

\end{document}